%% =====================================================================
%%  Capítulo 7 -- Controladores PID: nociones básicas
%%  Migrado al formato nuevo (noviembre de 2026) desde Handout6a_Ctrl.tex.
%%  Correcciones de fondo (ver MIGRACION.md): conversión ideal -> serie
%%  (tau_d^S usa tau_i^I, no tau_d^I); expansión de la forma serie; la forma
%%  paralelo no es "la única que permite quitar la integral"; historia
%%  (Taylor: pre-act fijo en 1936, ajustable en el Fulscope de 1939); se
%%  completaron las figuras y simulaciones marcadas "AÑADIR FIGURA";
%%  notación tau_i, tau_d -> T_i, T_d (convenciones del curso); MP -> %OS.
%% =====================================================================

\capitulo[cap:pn]{Controladores PID: nociones básicas}
  {Nunca, ni en las situaciones más desesperadas, he perdido mi sentido del humor, y soy de los que piensan que esta vida es de risa.}
  {Enrique Vila-Matas}

\begin{objetivos}[partes/07-pid-nociones.html]
  \item Explicar qué hace cada una de las tres acciones (proporcional, integral y derivativa) y cómo afecta la respuesta transitoria, el error y la estabilidad.
  \item Pasar los parámetros de un PID entre las formas ideal, paralelo y serie, y saber cuándo la conversión es posible.
  \item Reconocer los problemas prácticos del PID ---\emph{kicks}, ruido, \emph{wind-up}--- y las modificaciones que los resuelven: filtro derivativo, estructuras I-P e I-PD, anti-windup y ponderación de la referencia.
  \item Escribir el algoritmo digital del PID en forma posicional e incremental.
\end{objetivos}

%% ---------------------------------------------------------------------
\section{El controlador de tres términos}
\label{sec:pn-intro}
%% ---------------------------------------------------------------------

Hasta aquí hemos modelado, analizado la respuesta transitoria, la estabilidad y el error. Ahora empezamos a \emph{diseñar}, y lo hacemos con el controlador más usado en la industria: el de tres términos, \clave{proporcional, integral y derivativo} (PID). Más del noventa por ciento de los lazos de control de procesos son PI o PID~\citep{ACL05}. Los primeros fueron neumáticos, luego electrónicos y, desde los años setenta, digitales; la idea, sin embargo, es la misma desde hace un siglo.

\begin{historia}[De los timoneles a la Taylor]
En 1910 Elmer Sperry fundó la Sperry Gyroscope Company; su piloto automático para barcos, que acoplaba un girocompás al timón, hacía en la práctica lo que hacían los timoneles expertos, y resultó ser un PID mecánico~\citep{AsKu14}. En 1922 Nicolas Minorsky, que había observado a esos timoneles, propuso por escrito la ley de tres términos y analizó la estabilidad del lazo con un modelo lineal (capítulo~1)~\citep{Min22,Ben01}. En la industria de procesos el PID llegó con los controladores neumáticos. La Foxboro vendía desde 1931 el \emph{Stabilog}, con acción proporcional e integral (\emph{reset}), y la Taylor agregó hacia 1936 una acción derivativa fija (\emph{pre-act}) a su controlador. En 1939 aparecieron los dos primeros PID comerciales con las tres acciones ajustables: el \emph{Fulscope} rediseñado de la Taylor y el \emph{Stabilog} con \emph{hyper-reset} de la Foxboro~\citep{Ben84,Ben93}. Entre 1925 y 1935 se vendieron en Estados Unidos unos $75\,000$ controladores, y había unas $600$ empresas de instrumentos. Quedaba un problema: cómo escoger los parámetros. Lo resolvieron en 1942 John Ziegler y Nathaniel Nichols, de la Taylor~\citep{ZiNi42}, y en 1953 Cohen y Coon, también de la Taylor~\citep{CoCo53}; Geraldine Coon comprobó sus reglas en el computador analógico de la empresa. En 1959 se cerró el primer lazo industrial completamente digital, en la unidad de polimerización de la refinería de Texaco en Port Arthur (Texas).
\end{historia}

\begin{figure}[ht]
\centering
\includegraphics[width=0.9\linewidth]{cap07/lazo_pid}
\caption{Lazo de control con filtro de la referencia $G_f$, perturbación de carga $D$ y ruido de medida $N$.}
\label{fig:pn-lazo}
\end{figure}

En el lazo de la \cref{fig:pn-lazo}, $R$ es la referencia (\emph{set point}, SP), $Y$ la variable del proceso (PV), $U_c$ la salida del controlador (OC), $D$ una perturbación de carga y $N$ el ruido de medida. Mientras no digamos otra cosa, $G_f=1$ y $D=N=0$, y el error es $E=R-Y$. El \clave{PID ideal} es
\begin{equation}
  \frac{U_c(s)}{E(s)}=\Kc\left(1+\frac{1}{\Ti s}+\Td s\right),
  \qquad
  u_c(t)=\Kc\left(e(t)+\frac{1}{\Ti}\int_0^t e(\tau)\,\dif\tau+\Td\,\deriv{e}{t}\right),
  \label{eq:pn-ideal}
\end{equation}
con $\Kc$ la ganancia proporcional, $\Ti$ el \clave{tiempo integral} y $\Td$ el \clave{tiempo derivativo}.

\begin{ideaclave}
\textbf{Pasado, presente y futuro.}\enspace La acción proporcional mira el error \emph{ahora}; la integral, el área acumulada, es decir, el \emph{pasado}; la derivativa, la tendencia, es decir, hacia dónde va el error en el \emph{futuro} cercano (\cref{fig:pn-tiempos}).
\end{ideaclave}

\begin{figure}[ht]
\centering
\includegraphics[width=0.8\linewidth]{cap07/f_tiempos}
\caption{Las tres acciones del PID sobre la misma señal de error.}
\label{fig:pn-tiempos}
\end{figure}

%% ---------------------------------------------------------------------
\section{El controlador más sencillo: on/off}
\label{sec:pn-onoff}
%% ---------------------------------------------------------------------

Para entender el PID conviene empezar por el controlador más simple, el \clave{on/off} (o de dos posiciones):
\[
  u_c(t)=\begin{cases} u_{\max}, & e(t)>0,\\ u_{\min}, & e(t)<0. \end{cases}
\]
Solo mira el signo del error. Supone que el proceso tiene ganancia positiva: más salida del controlador, más PV.

\begin{ejemplo}[label=ej:pn-onoff]{Un horno on/off}
Un proceso $e^{-2s}/(10s+1)$ se controla con $r=50\%$, $u_{\max}=90\%$ y $u_{\min}=10\%$ (\cref{fig:pn-onoff}).
\solucion La PV oscila sin asentarse alrededor de la referencia, entre unos $42$ y $58\%$, y la salida del controlador salta entre sus dos valores cada pocos segundos. La oscilación la sostiene el tiempo muerto: cuando el controlador cambia, el proceso tarda $2$\,s en enterarse.
\end{ejemplo}

\begin{figure}[ht]
\centering
\includegraphics[width=0.85\linewidth]{cap07/f_onoff}
\caption{Ejemplo~\ref{ej:pn-onoff}. Arriba, referencia y PV; abajo, salida del controlador y error.}
\label{fig:pn-onoff}
\end{figure}

El on/off es barato, no hay nada que sintonizar y basta un sensor que distinga ``por encima'' de ``por debajo''. Por eso está en los termostatos de hornos, neveras y calentadores. Sus defectos son la oscilación permanente de la PV y los cambios bruscos de la salida, que desgastan el elemento final de control (una válvula, una bomba, un relé). En la práctica se le agrega una histéresis para que no conmute con el ruido.

\codigorepo{cap07_pid_nociones/C7_Ej1_OnOff.m}{onoff}{ejemplo~\ref{ej:pn-onoff}: el control on/off y su oscilación}

%% ---------------------------------------------------------------------
\section{Las tres acciones, una por una}
\label{sec:pn-acciones}
%% ---------------------------------------------------------------------

\subsection{Acción proporcional}

El on/off funciona razonablemente cuando el error es grande: si la PV está lejos, tiene sentido abrir la válvula del todo. El problema aparece cerca de la referencia. La solución natural es que la acción sea \emph{proporcional} al error dentro de una banda:
\[
  u_c(t)=\begin{cases} u_{\max}, & e(t)>e_0,\\ u_0+\Kc\,e(t), & -e_0\le e(t)\le e_0,\\ u_{\min}, & e(t)<-e_0. \end{cases}
\]
$u_0$ es la salida cuando el error es cero (el \emph{bias} o \emph{reset} manual) y $\Kc$ la ganancia proporcional (\cref{fig:pn-banda}). Dentro de la banda, $\Gc(s)=\Kc$.

\begin{figure}[ht]
\centering
\includegraphics[width=0.85\linewidth]{cap07/f_banda}
\caption{Característica estática del on/off y del proporcional con banda.}
\label{fig:pn-banda}
\end{figure}

\paragraph{El error en estado estacionario.} Despejando, $e(t)=\bigl(u_c(t)-u_0\bigr)/\Kc$: el error solo es cero si, en estado estacionario, la salida del controlador coincide con el \emph{bias}. Para un proceso de tipo~0 eso ocurre para una sola referencia. Con $\Gp=\Kproc/(\tau s+1)$ y $u_0=0$, el error ante escalón es $\ess=1/(1+\Kc\Kproc)$ (capítulo~6): baja al subir $\Kc$, pero no se anula, y una ganancia muy alta puede desestabilizar el lazo (capítulo~5).

\begin{cuidado}[Ajustar el \emph{bias} a mano no es robusto]
Con el proceso conocido, el \emph{bias} que anula el error es $u_0=r/\Gp(0)$. Pero basta que la ganancia del proceso cambie un poco (ensuciamiento, otro punto de operación) para que el error reaparezca. Lo que hace falta es un mecanismo que ajuste $u_0$ \emph{solo}: esa es la acción integral.
\end{cuidado}

\paragraph{Banda proporcional.} Muchos controladores industriales no muestran $\Kc$ sino la \clave{banda proporcional}: el porcentaje del rango del error que lleva la salida de un extremo al otro. Con $e_R$ y $u_R$ los rangos del error y de la salida,
\[
  PB\%=\frac{1}{\Kc}\,\frac{u_R}{e_R}\times100\%,
  \qquad\text{y con señales normalizadas al }100\%:\qquad PB\%=\frac{100}{\Kc}.
\]
Una banda estrecha es una ganancia grande. (Algunos fabricantes, como ABB, usan la primera forma; otros, como Foxboro, la segunda~\citep{ACL05}.)

\begin{pausa}
Un controlador tiene $PB=25\%$ con señales normalizadas. ¿Cuánto vale $\Kc$? ¿Y si el error se mide sobre un rango de $200\,\grados$C y la salida sobre $100\%$?
\respuesta{$\Kc=100/25=4$ (adimensional). Con los rangos: $\Kc=\frac{1}{0.25}\cdot\frac{100}{200}=2$\,\%/$\grados$C.}
\end{pausa}

\subsection{Acción integral}

La acción integral es proporcional al área acumulada del error:
\[
  u_c(t)=K_i\int_0^t e(\tau)\,\dif\tau,\qquad K_i=\frac{\Kc}{\Ti},\qquad \Gc(s)=\frac{K_i}{s}.
\]
Mientras haya error, la salida sigue cambiando; solo se detiene cuando el error es cero. Ese polo en el origen sube en uno el \clave{tipo} del sistema: el error ante escalón, en la referencia o en una perturbación de carga, se anula.

\begin{ideaclave}
\textbf{La integral ajusta el \emph{bias} sola.}\enspace En un PI, $u_0$ se reemplaza por $\frac{\Kc}{\Ti}\int_0^t e\,\dif\tau$, que crece o decrece hasta encontrar el valor que deja el error en cero, sin conocer el proceso. Así se construían los PI neumáticos, con un lazo de \emph{reset automático} (\cref{fig:pn-reset}): $U_c=\Kc E+\frac{1}{\Ti s+1}U_c$, es decir, $\frac{U_c}{E}=\Kc\Bigl(1+\frac{1}{\Ti s}\Bigr)$.
\end{ideaclave}

\begin{figure}[ht]
\centering
\includegraphics[width=0.55\linewidth]{cap07/pi_reset}
\caption{PI con reajuste automático (\emph{automatic reset}).}
\label{fig:pn-reset}
\end{figure}

La acción integral tiene memoria: cuando el error llega a cero, la integral no es cero, porque guarda todo lo anterior. La \cref{fig:pn-integral} (izquierda) lo muestra con $\Gc=1/s$ y $\Gp=1/(s+1)$: la salida del controlador sobrepasa su valor final y solo empieza a bajar cuando el error se vuelve negativo, es decir, después de que la salida ya pasó la referencia. De ahí viene el sobrepico. A la derecha, al reducir $\Ti$ (más acción integral) la respuesta es más rápida pero menos amortiguada: el lazo es $s^2+s+1/\Ti$, con $\zeta=\tfrac12\sqrt{\Ti}$.

\begin{figure}[ht]
\centering
\includegraphics[width=\linewidth]{cap07/f_integral}
\caption{Izquierda: control integral puro; la salida del controlador ($u$) sigue subiendo mientras el error es positivo. Derecha: menos $\Ti$, más rápido y más oscilatorio.}
\label{fig:pn-integral}
\end{figure}

El PI resuelve los dos problemas anteriores: la oscilación del on/off y el error del proporcional. Por eso es, de lejos, la estructura más usada en la industria; si no hay exigencias fuertes de velocidad, la acción derivativa sobra. Su precio es una pérdida de estabilidad (un polo más en el origen) y un fenómeno que veremos en la \cref{sec:pn-modificaciones}: el \emph{wind-up}. Algunos fabricantes usan, en lugar de $\Ti$, la \clave{tasa de reajuste} $1/\Ti$, en repeticiones por minuto: cuántas veces por minuto la integral repite la acción proporcional ante un error constante.

\subsection{Acción derivativa}

El PI mira el pasado y el presente, pero no se anticipa. En la \cref{fig:pn-derivada} los dos errores tienen el mismo valor en $t_1$ y la misma área hasta $t_1$, así que un PI daría exactamente la misma salida. Sin embargo, en (a) el error cae rápido y el controlador debería frenar para no pasarse, y en (b) el error está volviendo a subir y debería reaccionar ya.

\begin{figure}[ht]
\centering
\includegraphics[width=0.8\linewidth]{cap07/f_derivada}
\caption{Dos errores con el mismo valor y la misma área en $t_1$: un PI no los distingue; la derivada sí.}
\label{fig:pn-derivada}
\end{figure}

La acción derivativa es proporcional a la pendiente del error: $u_c=K_d\,\dot e$, con $K_d=\Kc\Td$. Casi nunca se usa sola, sino con la proporcional, y entonces tiene una interpretación muy clara. Por Taylor, $e(t+\Td)\approx e(t)+\Td\,\dot e(t)$, así que
\[
  \Kc\bigl(e(t)+\Td\,\dot e(t)\bigr)\approx\Kc\,e(t+\Td):
\]
\begin{ideaclave}
\textbf{El PD es un proporcional que mira $\Td$ segundos adelante.}\enspace La salida en $t$ se basa en el error que se \emph{predice} para $t+\Td$, extrapolando en línea recta. Por eso los fabricantes la llamaron \emph{pre-act} o \emph{rate action}.
\end{ideaclave}

La derivada sirve cuando el lazo se beneficia de anticiparse, como en procesos de temperatura de orden alto. Pero tiene dos problemas. $\Gc=K_d s$ es una función de transferencia \alerta{impropia} (más ceros que polos): no se puede realizar físicamente y amplifica sin límite las altas frecuencias, es decir, el ruido. Y reacciona con violencia a los saltos de la referencia. Las dos cosas se resuelven en la \cref{sec:pn-modificaciones}.

%% ---------------------------------------------------------------------
\section{Tres formas de escribir un PID}
\label{sec:pn-estructuras}
%% ---------------------------------------------------------------------

Los controladores comerciales no implementan todos el mismo algoritmo, y los parámetros que da un método de sintonización dependen de la forma. Las tres clásicas son:
\begin{align}
  \text{ideal (ISA, estándar):}\qquad & \Gc(s)=\Kc^I\left(1+\frac{1}{\Ti^I s}+\Td^I s\right), \label{eq:pn-forma-ideal}\\
  \text{paralelo:}\qquad & \Gc(s)=\Kc^P+\frac{K_i^P}{s}+K_d^P s, \label{eq:pn-forma-paralelo}\\
  \text{serie (interactiva):}\qquad & \Gc(s)=\Kc^S\left(1+\frac{1}{\Ti^S s}\right)\left(1+\Td^S s\right). \label{eq:pn-forma-serie}
\end{align}
Los métodos de sintonización del capítulo~8 dan los parámetros en forma ideal; si el controlador es de otra forma, hay que convertirlos.

\paragraph{Paralelo e ideal.} Las tres acciones son independientes en las dos, y la conversión es inmediata:
\[
  \Kc^I=\Kc^P,\qquad \Ti^I=\frac{\Kc^P}{K_i^P},\qquad \Td^I=\frac{K_d^P}{\Kc^P}.
\]
La forma paralelo es la más general: es la única que admite $\Kc^P=0$, es decir, un controlador integral o integral-derivativo puro. (Quitar la integral o la derivada se puede en cualquiera de las tres, con $\Ti\to\infty$ o $\Td=0$.) La usan, entre otros, Siemens, Modicon y Allen-Bradley.

\paragraph{Serie e ideal.} Multiplicando la forma serie,
\[
  \Kc^S\left(1+\frac{1}{\Ti^S s}\right)\left(1+\Td^S s\right)
  =\Kc^S\left(1+\frac{\Td^S}{\Ti^S}+\frac{1}{\Ti^S s}+\Td^S s\right),
\]
y agrupando como en la ideal,
\begin{equation}
  \Kc^I=\Kc^S\left(1+\frac{\Td^S}{\Ti^S}\right),\qquad
  \Ti^I=\Ti^S+\Td^S,\qquad
  \Td^I=\frac{\Ti^S\Td^S}{\Ti^S+\Td^S}.
  \label{eq:pn-serie-ideal}
\end{equation}
Las acciones interactúan: cambiar $\Td^S$ cambia la ganancia y el tiempo integral equivalentes. El camino inverso no siempre existe. Imponiendo \eqref{eq:pn-serie-ideal}, $\Ti^S+\Td^S=\Ti^I$ y $\Ti^S\Td^S=\Ti^I\Td^I$, así que $\Ti^S$ y $\Td^S$ son las raíces de $z^2-\Ti^Iz+\Ti^I\Td^I=0$:
\begin{equation}
  \Kc^S=\frac{\Kc^I}{2}\bigl(1+\rho\bigr),\qquad
  \Ti^S=\frac{\Ti^I}{2}\bigl(1+\rho\bigr),\qquad
  \Td^S=\frac{\Ti^I}{2}\bigl(1-\rho\bigr),\qquad
  \rho=\sqrt{1-4\,\frac{\Td^I}{\Ti^I}} ,
  \label{eq:pn-ideal-serie}
\end{equation}
que solo son reales si \alerta{$\Ti^I\ge4\Td^I$}.

\begin{ideaclave}
\textbf{Los ceros del PID.}\enspace Un PID tiene un polo en el origen y dos ceros, las raíces de $\Ti\Td s^2+\Ti s+1$. La forma serie solo puede tener ceros reales (en $-1/\Ti^S$ y $-1/\Td^S$). La ideal admite ceros complejos, cuando $\Ti<4\Td$: es más general. En $\Ti=4\Td$ los dos ceros coinciden.
\end{ideaclave}

La forma serie fue la primera que se construyó con tecnología neumática, y muchos fabricantes la conservaron por tradición y porque se dice que es más fácil de sintonizar a mano.

\begin{ejemplo}[label=ej:pn-conversion]{De serie a ideal, y de vuelta}
Un controlador en forma serie tiene $\Kc^S=2$, $\Ti^S=8$\,s y $\Td^S=2$\,s. Halle la forma ideal y regrese.
\solucion Con \eqref{eq:pn-serie-ideal}: $\Kc^I=2(1+2/8)=2.5$, $\Ti^I=10$\,s y $\Td^I=16/10=1.6$\,s. De vuelta, $\rho=\sqrt{1-4\cdot1.6/10}=0.6$, y \eqref{eq:pn-ideal-serie} da $\Kc^S=2.5\cdot0.8=2$, $\Ti^S=5\cdot1.6=8$\,s y $\Td^S=5\cdot0.4=2$\,s.
\end{ejemplo}

\codigorepo{cap07_pid_nociones/C7_Ej3_Estructuras.m}{conversion}{ejemplo~\ref{ej:pn-conversion} con la función \cmd{convertirPID}}

\begin{pausa}
¿Tiene forma serie el PID ideal con $\Kc=5$, $\Ti=4$\,s y $\Td=2$\,s?
\respuesta{No: $4\Td=8>\Ti=4$. Los ceros, raíces de $8s^2+4s+1$, son $-0.25\pm j0.25$.}
\end{pausa}

%% ---------------------------------------------------------------------
\section{Un ejemplo para verlas actuar}
\label{sec:pn-ejemplo}
%% ---------------------------------------------------------------------

\begin{ejemplo}[label=ej:pn-mra]{Masa, resorte y amortiguador}
El sistema $M\ddot x+b\dot x+kx=F$, con $M=1$\,kg, $b=10$\,N\,s/m y $k=20$\,N/m, tiene
\[
  \Gp(s)=\frac{X(s)}{F(s)}=\frac{1}{s^2+10s+20}.
\]
Se quiere, con retroalimentación unitaria y un PID en forma paralelo, un tiempo de subida corto, el menor sobrepico posible y error cero ante escalón.
\solucion \textbf{Lazo abierto.} Los polos son $-2.76$ y $-7.24$ ($\zeta=1.12$): sin sobrepico, con $\tr=0.88$\,s, $\ts=1.60$\,s y valor final $1/20$. Respecto a una referencia unitaria, el error es del $95\%$.

\textbf{P.} El lazo es $s^2+10s+(20+\Kc)$: $\wn=\sqrt{20+\Kc}$, $\zeta=5/\sqrt{20+\Kc}$ y $\ess=20/(20+\Kc)$. Subir $\Kc$ reduce el tiempo de subida y el error, pero baja $\zeta$ y sube el sobrepico. Con $\Kc=300$: $\ess=0.0625$, $\zeta=0.28$ y $\OS=40\%$ (\cref{fig:pn-mra}, arriba a la izquierda).

\textbf{PD.} Con $\Kc=300$, la derivada agrega amortiguamiento: el lazo es $s^2+(10+K_d)s+320$. Con $K_d=20$, $\zeta=0.84$ y el sobrepico cae al $5.8\%$; el error no cambia ($K_ds$ se anula en $s=0$). El PD introduce además un cero en $-\Kc/K_d=-15$; por eso el sobrepico real es mayor que el que predice el segundo orden sin cero ($0.8\%$).

\textbf{PI.} Con $\Kc=300$, la integral anula el error, pero el lazo es de tercer orden, $s^3+10s^2+320s+K_i$, estable solo si $K_i<3200$ (Routh). Con $K_i=200$: $\ess=0$ y $\OS=36\%$.

\textbf{PID.} Con $\Kc=350$, $K_i=300$ y $K_d=50$, la respuesta no tiene sobrepico, sube en $0.05$\,s, se establece en $0.83$\,s y no tiene error. Los ceros del controlador ($-1$ y $-6$) casi cancelan los polos más lentos del lazo cerrado.
\end{ejemplo}

\begin{figure}[ht]
\centering
\includegraphics[width=\linewidth]{cap07/f_mra}
\caption{Ejemplo~\ref{ej:pn-mra}: efecto de cada acción sobre la respuesta al escalón.}
\label{fig:pn-mra}
\end{figure}

\codigorepo{cap07_pid_nociones/C7_Ej2_Acciones.m}{pid}{ejemplo~\ref{ej:pn-mra}: P, PD, PI y PID sobre el sistema masa-resorte-amortiguador}

\begin{table}[ht]
\centering
\caption{Efecto típico de aumentar cada parámetro de un PID sobre una planta estable~\citep{ACL05}. Son tendencias, no leyes: el problema~\ref{prob:pn-mra} pide comprobarlas.}
\label{tab:pn-efectos}
\small
\begin{tabular}{@{}lccccc@{}}
\toprule
Aumentar & $\tr$ & $\OS$ & $\ts$ & $\ess$ & Estabilidad\\
\midrule
$\Kc$ & baja & sube & sube un poco & baja & empeora\\
$K_i=\Kc/\Ti$ & baja un poco & sube & sube & se anula & empeora\\
$K_d=\Kc\Td$ & baja un poco & baja & baja & casi no cambia & mejora (si no hay ruido)\\
\bottomrule
\end{tabular}
\end{table}

\begin{receta}[Ajuste manual de un PID]
\begin{enumerate}
  \item Mire la respuesta en lazo abierto y decida qué hay que mejorar.
  \item Agregue acción proporcional para acelerar la respuesta.
  \item Agregue acción derivativa si el sobrepico es excesivo.
  \item Agregue acción integral para eliminar el error en estado estacionario.
  \item Ajuste las tres hasta lograr la respuesta deseada, y verifique la estabilidad.
  \item No use las tres si no hace falta: la derivada trae problemas (\cref{sec:pn-modificaciones}). Mantenga el controlador tan simple como se pueda.
\end{enumerate}
\end{receta}

En resumen: el \textbf{P} acelera pero deja error (salvo en procesos integradores); el \textbf{PI} elimina el error pero oscila más y pierde estabilidad al subir $\Kc$; el \textbf{PD} anticipa y amortigua, pero deja el error del P; el \textbf{PID} junta las ventajas, a costa de amplificar el ruido. En el capítulo~10 (\cref{sec:rf-pid}) veremos lo mismo en frecuencia: un PI es un compensador de atraso, un PD uno de adelanto y un PID uno de atraso-adelanto.

%% ---------------------------------------------------------------------
\section{Del PID de libro al PID de planta}
\label{sec:pn-modificaciones}
%% ---------------------------------------------------------------------

El PID ideal tiene problemas que aparecen apenas se conecta a un proceso real: saltos en la salida cuando cambia la referencia, ruido amplificado y un comportamiento extraño cuando el actuador se satura. Casi todos los controladores comerciales traen opciones para resolverlos; veamos cuáles y por qué funcionan.

\subsection{El \emph{proportional kick} y la estructura I-P}

Si el lazo está en equilibrio ($e\approx0$) y la referencia salta, el error salta con ella, y la acción proporcional lleva ese salto completo a la salida del controlador: un golpe al actuador que se llama \clave{proportional kick}. La solución es que el término proporcional actúe sobre la salida del proceso y no sobre el error, mientras la integral sigue sobre el error (estructura \clave{I-P}, \cref{fig:pn-ip}):
\[
  U_c(s)=\frac{\Kc}{\Ti s}E(s)-\Kc Y(s).
\]
La salida del controlador ya no salta, y el rechazo de perturbaciones es el mismo (para $R=0$ las dos estructuras coinciden). A cambio, la respuesta a la referencia es más lenta.

\begin{figure}[ht]
\centering
\includegraphics[width=0.72\linewidth]{cap07/estructura_ip}
\caption{Estructura I-P.}
\label{fig:pn-ip}
\end{figure}

\subsection{La derivada y el ruido: el filtro derivativo}

Un ruido de medida $n(t)=A\sin\omega t$ que pase por la derivada pura produce $u_c=AK_d\,\omega\cos\omega t$: \alerta{la amplitud crece con la frecuencia}. Con $\Kc=3$, $\Td=2$ y $A=0.01$, un ruido de $100$\,rad/s aparece en la salida del controlador con amplitud $6$. La solución es limitar la ganancia de alta frecuencia con un \clave{filtro derivativo}:
\begin{equation}
  \Gc(s)=\Kc\left(1+\frac{1}{\Ti s}+\frac{\Td s}{\dfrac{\Td}{N}s+1}\right).
  \label{eq:pn-filtro}
\end{equation}
Para $\omega\ll N/\Td$ la derivada funciona como antes; por encima, su ganancia se aplana en $\Kc N$, y la del controlador en $\Kc(1+N)$ (\cref{fig:pn-ruido}). El controlador queda propio (grado relativo cero) y realizable. $N$ suele estar entre $8$ y $16$~\citep{ACL05}; otros textos dan $5\le N\le20$~\citep{DoBi08}. Con $N=10$ el ruido del ejemplo baja de $6$ a $0.30$.

\begin{figure}[ht]
\centering
\includegraphics[width=0.8\linewidth]{cap07/f_ruido}
\caption{Magnitud del PID ideal ($\Kc=3$, $\Ti=8$, $\Td=2$) sin filtro y con filtro derivativo. Las líneas punteadas son las cotas $\Kc(1+N)$.}
\label{fig:pn-ruido}
\end{figure}

\begin{profundiza}[Otros filtros]
Hay variantes: un filtro de segundo orden en la derivada, $\Td s/\bigl(1+\frac{\Td s}{N}+\frac12(\frac{\Td s}{N})^2\bigr)$~\citep{gerry2004pid}; la forma serie con filtro, $\Kc\bigl(1+\frac{1}{\Ti s}\bigr)\frac{\Td s+1}{\Td s/N+1}$; o un filtro sobre toda la salida, $\Kc\bigl(1+\frac{1}{\Ti s}+\Td s\bigr)\frac{1}{(\tau_fs+1)^k}$ con $k=1$ o $2$. El diseño del filtro es un compromiso entre atenuar el ruido y no perder la acción derivativa útil.
\end{profundiza}

\subsection{El \emph{derivative kick} y la estructura I-PD}

Aun con filtro, si la derivada actúa sobre el error, un escalón en la referencia produce un pico enorme: la derivada (filtrada) de un escalón es un pulso de altura $\Kc\Td/(\Td/N)=\Kc N$. Es el \clave{derivative kick}. Con el PID $350/300/50$ del ejemplo~\ref{ej:pn-mra} y un filtro de constante $0.01$\,s, el pico es de unas $5300$ unidades (\cref{fig:pn-kicks}). Como la referencia es constante a trozos, su derivada no aporta nada útil entre saltos; por eso se deriva solo la salida. Con la proporcional también sobre la salida se obtiene la estructura \clave{I-PD} (\cref{fig:pn-ipd}).

\begin{figure}[ht]
\centering
\includegraphics[width=0.72\linewidth]{cap07/estructura_ipd}
\caption{Estructura I-PD: solo la integral ve el error.}
\label{fig:pn-ipd}
\end{figure}

\begin{figure}[ht]
\centering
\includegraphics[width=\linewidth]{cap07/f_kicks}
\caption{PID $350/300/50$ sobre $1/(s^2+10s+20)$. Estándar: pico de unas $5300$ unidades en la salida del controlador. PI-D (derivada sobre la salida): queda el salto proporcional de $350$. I-PD: sin saltos, pero más lento.}
\label{fig:pn-kicks}
\end{figure}

\codigorepo{cap07_pid_nociones/C7_Ej5_Kicks.m}{kicks}{las tres estructuras ante un escalón de referencia}

\subsection{\emph{Wind-up} y anti-windup}

Todo actuador tiene límites: una válvula no abre más del $100\%$. Si un cambio grande de referencia lleva la salida del controlador al límite, el lazo queda, en la práctica, abierto: el actuador ya no responde a lo que pida el controlador. Pero el error sigue siendo grande, y la integral sigue acumulándolo (\emph{winds up}). Cuando la salida del proceso por fin alcanza la referencia, la integral está inflada y mantiene el actuador saturado hasta que un error negativo igual de grande la descarga: sobrepico grande y establecimiento lento. Es el \clave{wind-up} del integrador.

\begin{ejemplo}[label=ej:pn-windup]{\emph{Wind-up} en un proceso con tiempo muerto}
Un PID ideal con $\Kc=3$, $\Ti=8$\,s, $\Td=2$\,s (y $N=10$) controla $\Gp(s)=e^{-4s}/(10s+1)$, con el actuador limitado a $[0,\,1.5]$~\citep{Vis06}.
\solucion Ante un escalón unitario, el actuador queda saturado $25$\,s en total; el término integral llega a más de $3$ cuando el actuador solo puede dar $1.5$. Resultado: $\OS=38\%$ y $\ts=45$\,s (\cref{fig:pn-windup}, en rojo).
\end{ejemplo}

Las correcciones \clave{anti-windup} más usadas son dos.
\begin{itemize}
  \item \textbf{Integración condicional} (\emph{clamping}): se detiene la integral mientras el actuador esté saturado y el error empuje hacia el mismo lado ($u_c$ saturado y $(v-u)\,e>0$). No hay que sintonizar nada. En el ejemplo: $\OS=3.7\%$, $\ts=29$\,s.
  \item \textbf{Recálculo} (\emph{back-calculation}, \cref{fig:pn-antiwindup}): se retroalimenta a la integral la diferencia entre la salida saturada $u$ y la que pedía el controlador $v$,
  \[
    \frac{\dif I}{\dif t}=\frac{\Kc}{\Ti}\,e+\frac{1}{T_t}\,(u-v).
  \]
  Mientras no hay saturación, $u=v$ y no pasa nada; con saturación, la integral se descarga con constante de tiempo $T_t$. Funciona como un observador: estima el estado correcto del controlador. Se recomienda $T_t=\sqrt{\Ti\Td}$~\citep{astrom1995pid}, o $T_t=\Ti$ para un PI~\citep{atherton1995pid}. En el ejemplo, con $T_t=4$: $\OS=3.3\%$, $\ts=30$\,s; con $T_t=8$: $\OS=11\%$.
\end{itemize}

\begin{figure}[ht]
\centering
\includegraphics[width=0.85\linewidth]{cap07/antiwindup}
\caption{PID con anti-windup por recálculo (\emph{back-calculation}).}
\label{fig:pn-antiwindup}
\end{figure}

\begin{figure}[ht]
\centering
\includegraphics[width=\linewidth]{cap07/f_windup}
\caption{Ejemplo~\ref{ej:pn-windup}. Sin corrección, la integral (punteada) sube mucho más de lo que el actuador puede dar; con anti-windup se mantiene acotada.}
\label{fig:pn-windup}
\end{figure}

\codigorepo{cap07_pid_nociones/C7_Ej6_Windup.m}{windup}{ejemplo~\ref{ej:pn-windup}: sin corrección, con recálculo y con integración condicional}

\begin{modelo}[AntiWindup.slx]
El script \cmd{crearModelo\_AntiWindup.m} arma el ejemplo con el bloque \emph{PID Controller} de \simulink, que trae la saturación y el anti-windup incorporados (parámetros \emph{Limit output} y \emph{Anti-windup method}). Compare las dos copias del lazo y cambie $T_t$.
\end{modelo}

\begin{cuidado}[Dónde aparece el \emph{wind-up}]
El \emph{wind-up} es más grave ante cambios grandes de la referencia que ante perturbaciones de carga~\citep{VPPH97}, y en procesos de orden bajo con tiempo muerto pequeño respecto a la constante de tiempo, donde el controlador puede ser agresivo. Un controlador que nunca satura no lo necesita, pero en planta casi todos saturan alguna vez (arranques, cambios de producto).
\end{cuidado}

\subsection{Ponderación de la referencia: dos grados de libertad}

Un lazo sintonizado para rechazar perturbaciones rápido suele responder con mucho sobrepico a los cambios de referencia. Con un solo controlador sobre el error no se pueden tener las dos cosas. La salida es separar cómo actúa el controlador sobre la referencia y sobre la salida (\cref{fig:pn-2gdl}):
\[
  u_c=\Kc\Bigl(\beta\,r-y\Bigr)+\frac{\Kc}{\Ti}\int(r-y)\,\dif\tau+\Kc\Td\,\deriv{}{t}\bigl(\gamma\,r-y\bigr),
  \qquad 0\le\beta,\gamma\le1 .
\]
La integral siempre actúa sobre el error completo, para no perder el error cero. Con $\beta=\gamma=1$ es el PID estándar; con $\beta=1$, $\gamma=0$, el PI-D; con $\beta=\gamma=0$, el I-PD. Esta es la estructura del \clave{PID ISA} (o controlador $\beta$-$\gamma$). Equivale a un filtro sobre la referencia,
\[
  G_f(s)=\frac{1+\bigl(\beta\Ti+\frac{\Td}{N}\bigr)s+\Ti\Td\bigl(\gamma+\frac{\beta}{N}\bigr)s^2}{1+\bigl(\Ti+\frac{\Td}{N}\bigr)s+\Ti\Td\bigl(1+\frac{1}{N}\bigr)s^2},
\]
que \alerta{no aparece en la función de transferencia de la perturbación a la salida}: el rechazo de perturbaciones se sintoniza con $\Kc$, $\Ti$, $\Td$, y la respuesta a la referencia, después, con $\beta$ y $\gamma$.

\begin{figure}[ht]
\centering
\includegraphics[width=0.62\linewidth]{cap07/dos_gdl}
\caption{Controlador de dos grados de libertad: $G_{sp}=\Kc\bigl(\beta+\frac{1}{\Ti s}+\gamma\Td s\bigr)$ sobre la referencia y $\Gc$ sobre la salida.}
\label{fig:pn-2gdl}
\end{figure}

\begin{ejemplo}[label=ej:pn-beta]{Un PI agresivo con ponderación}
Un PI con $\Kc=20$ y $\Ti=0.3$\,s controla $\Gp=1/\bigl((s+1)(s+5)\bigr)$; está sintonizado para rechazar perturbaciones (pico de $0.032$ ante un escalón unitario a la entrada de la planta, frente a $0.11$ de un PI conservador con $\Kc=2$, $\Ti=1$).
\solucion Ante la referencia, con $\beta=1$ el sobrepico es del $51\%$; con $\beta=0.5$, del $33\%$, y con $\beta=0$, del $27\%$ (\cref{fig:pn-beta}). La respuesta a la perturbación es la misma para todos. Note que ningún $\beta$ baja el sobrepico de $27\%$: con $\beta=0$ queda solo la integral sobre la referencia, y la oscilación la imponen los polos del lazo cerrado, que $\beta$ no mueve.
\end{ejemplo}

\begin{figure}[ht]
\centering
\includegraphics[width=\linewidth]{cap07/f_beta}
\caption{Ejemplo~\ref{ej:pn-beta}: $\beta$ cambia la respuesta a la referencia y no la respuesta a la perturbación.}
\label{fig:pn-beta}
\end{figure}

\codigorepo{cap07_pid_nociones/C7_Ej7_Ponderacion.m}{beta}{ejemplo~\ref{ej:pn-beta}: ponderación de la referencia}

%% ---------------------------------------------------------------------
\section{Implementación digital}
\label{sec:pn-digital}
%% ---------------------------------------------------------------------

Hoy casi todos los PID corren en un computador que muestrea el error cada $T_s$ segundos. Hay varias formas de discretizar~\citep{AsWi97}; la más sencilla aproxima la integral por una suma y la derivada por una diferencia hacia atrás:
\[
  \int_0^t e\,\dif\tau\approx T_s\sum_{i=1}^{k}e[i],\qquad \deriv{e}{t}\approx\frac{e[k]-e[k-1]}{T_s}.
\]
El resultado es el \clave{algoritmo posicional}, que calcula la salida completa en cada muestra:
\begin{equation}
  u_c[k]=\Kc\left(e[k]+\frac{T_s}{\Ti}\sum_{i=1}^{k}e[i]+\frac{\Td}{T_s}\bigl(e[k]-e[k-1]\bigr)\right).
  \label{eq:pn-posicional}
\end{equation}
Restando \eqref{eq:pn-posicional} en $k$ y en $k-1$ se obtiene el \clave{algoritmo incremental} (o de velocidad), que calcula solo el cambio:
\begin{equation}
  u_c[k]=u_c[k-1]+\Kc\left[\Bigl(1+\frac{T_s}{\Ti}+\frac{\Td}{T_s}\Bigr)e[k]-\Bigl(1+\frac{2\Td}{T_s}\Bigr)e[k-1]+\frac{\Td}{T_s}\,e[k-2]\right].
  \label{eq:pn-incremental}
\end{equation}

\begin{ideaclave}
\textbf{Por qué gusta el incremental.}\enspace No guarda la suma de todos los errores, se lleva bien con actuadores que reciben incrementos (motores de paso) y, si la salida se satura y se recorta el valor que se guarda como $u_c[k-1]$, no acumula integral: es naturalmente resistente al \emph{wind-up}. Note el factor $\Td/T_s$: con muestreo rápido amplifica mucho cualquier cambio entre dos muestras, incluido el ruido y la cuantización del conversor.
\end{ideaclave}

\begin{ejemplo}[label=ej:pn-digital]{Coeficientes del algoritmo incremental}
Para $\Kc=3$, $\Ti=8$\,s, $\Td=2$\,s y $T_s=0.1$\,s,
\[
  u_c[k]=u_c[k-1]+63.04\,e[k]-123\,e[k-1]+60\,e[k-2].
\]
Simulado con $\Gp=1/(10s+1)$, da la misma señal de control que el posicional (diferencias del orden de $10^{-13}$).
\end{ejemplo}

\codigorepo{cap07_pid_nociones/C7_Ej8_Digital.m}{digital}{ejemplo~\ref{ej:pn-digital}: posicional e incremental}

%% ---------------------------------------------------------------------
\section{Para llevar}
\label{sec:pn-cierre}
%% ---------------------------------------------------------------------

\begin{parallevar}
\begin{enumerate}
  \item PID ideal: $\Kc\bigl(1+\frac{1}{\Ti s}+\Td s\bigr)$. P mira el presente, I el pasado (área), D el futuro (tendencia): un PD es un proporcional que mira $\Td$ segundos adelante.
  \item El on/off oscila; el P deja error en procesos de tipo~0; la I lo elimina ajustando el \emph{bias} sola, a costa de estabilidad; la D amortigua, pero amplifica el ruido.
  \item Formas ideal, paralelo y serie. Serie $\to$ ideal siempre; ideal $\to$ serie solo si $\Ti\ge4\Td$. Los métodos de sintonización dan la forma ideal.
  \item Filtro derivativo con $N\approx8$--$16$: ganancia acotada en $\Kc(1+N)$ y controlador realizable.
  \item Los \emph{kicks} se evitan sacando la P y la D del camino de la referencia (I-P, PI-D, I-PD). La ponderación $\beta$, $\gamma$ separa la respuesta a la referencia del rechazo de perturbaciones.
  \item Con actuador saturado, la integral se infla (\emph{wind-up}). Anti-windup: integración condicional o recálculo con $T_t\approx\sqrt{\Ti\Td}$.
  \item Digital: algoritmo posicional o incremental; el incremental resiste mejor el \emph{wind-up}.
\end{enumerate}
\end{parallevar}

\begin{profundiza}[Correcciones respecto a la versión anterior del capítulo]
(i) La conversión ideal $\to$ serie daba $\Td^S=\frac{\Td^I}{2}(1-\rho)$; es $\frac{\Ti^I}{2}(1-\rho)$. (ii) La expansión de la forma serie tenía los términos mal agrupados. (iii) La forma paralelo no es ``la única que permite desactivar la integral'': es la única que admite $\Kc=0$. (iv) Historia: la Taylor introdujo una derivada fija hacia 1936; en 1939 salieron el Fulscope y el Stabilog con las tres acciones ajustables. Se retiró que la ganancia del Stabilog ``podía ajustarse entre 0.7 y 100'' y la mención al \emph{Trans Scan Log}, que no se encontraron en las fuentes. (v) Se completaron las figuras y simulaciones pendientes (on/off, kicks, \emph{wind-up}, ponderación) y se calcularon todos los números de los ejemplos. (vi) Notación: $\tau_i,\tau_d\to\Ti,\Td$; MP $\to\OS$; ``realimentar'' $\to$ retroalimentar; ``Allan Braedley'' $\to$ Allen-Bradley.
\end{profundiza}

%% ---------------------------------------------------------------------
\section{Problemas}
\label{sec:pn-problemas}
%% ---------------------------------------------------------------------

Los problemas marcados con \etiquetaM{} requieren \matlab/\simulink; los marcados con \etiquetaEstrella{} son de profundización.

\begin{problema}[formas del PID]\label{prob:pn-formas}
(a) Expanda la forma serie y deduzca \eqref{eq:pn-serie-ideal}. (b) Un controlador en forma serie tiene $\Kc^S=2$, $\Ti^S=8$\,s y $\Td^S=2$\,s; halle la forma ideal. (c) Muestre que $\Ti^S$ y $\Td^S$ son las raíces de $z^2-\Ti^Iz+\Ti^I\Td^I=0$, obtenga \eqref{eq:pn-ideal-serie} y verifique con (b). (d) Muestre, con el discriminante, que la conversión exige $\Ti^I\ge4\Td^I$. (e) Verifique que el PID ideal con $\Kc=5$, $\Ti=4$\,s y $\Td=2$\,s no tiene forma serie y calcule sus ceros. (f) Pase el PID paralelo $350/300/50$ del ejemplo~\ref{ej:pn-mra} a forma ideal y diga si tiene forma serie.
\end{problema}

\begin{problema}[las tendencias, con números]\label{prob:pn-mra}
Con el sistema del ejemplo~\ref{ej:pn-mra}: (a) en lazo abierto, calcule el valor final, $\zeta$, $\wn$, $\tr$ y $\ts$, y verifique con \cmd{stepinfo}; (b) con un P, halle $\ess$, $\zeta$, $\wn$ y $\OS$ en función de $\Kc$ y tabule para $\Kc=10$, $100$ y $300$; (c) con $\Kc=300$ y $K_d=10$, $20$, ubique el cero, explique por qué $\ess$ no cambia y compare el $\OS$ con el que predice el segundo orden sin el cero; (d) con $\Kc=300$, halle por Routh el rango de $K_i$ y evalúe $K_i=70$ y $200$; (e) \etiquetaM{} arme una tabla P, PD, PI, PID con $\tr$, $\OS$, $\ts$, $\ess$ y $|u_c|_{\max}$, y verifique una por una las tendencias de la \cref{tab:pn-efectos}. ¿Alguna no se cumple aquí?
\end{problema}

\begin{problema}[banda proporcional y \emph{bias}]\label{prob:pn-banda}
Un lazo de temperatura tiene un transmisor de $0$--$200\,\grados$C, una válvula de $0$--$100\%$ y un proceso de ganancia estática $\Kproc=2.5\,\grados$C/\%, con control proporcional. (a) Con $PB=25\%$, calcule $\Kc$ con las dos convenciones y explique por qué difieren. (b) Con la segunda y $u_0=50\%$, halle el error para $r=120\,\grados$C. (c) Halle el $u_0$ que lo anula. (d) Con ese $u_0$, ¿cuál es el error si la ganancia cae a $2.2$? ¿Y si sube a $2.8$? (e) Explique, con esos números, por qué el \emph{bias} manual no es robusto y qué hace la integral. (f) Con $u_0=50\%$ y $r=150\,\grados$C, ¿qué banda proporcional deja el error por debajo de $2\,\grados$C? ¿Es razonable?
\end{problema}

\begin{problema}[ruido y filtro derivativo]\label{prob:pn-ruido}
Con el PID ideal $\Kc=3$, $\Ti=8$\,s, $\Td=2$\,s: (a) para un ruido $0.01\sin\omega t$, calcule su amplitud en la salida del controlador con la derivada pura para $\omega=0.1$, $10$ y $100$\,rad/s; (b) muestre que la derivada filtrada tiene ganancia acotada en alta frecuencia y halle la cota; (c) halle la frecuencia a partir de la cual actúa el filtro y la amplitud del ruido en $\omega=100$ para $N=5$, $10$ y $20$; (d) \etiquetaM{} grafique $|\Gc(j\omega)|$ con y sin filtro; (e) verifique que el controlador filtrado es propio; (f) \etiquetaM{} simule el lazo con $\Gp=1/(10s+1)$ y ruido blanco en la medida, y mida la desviación estándar de $u_c$ para cada $N$.
\end{problema}

\begin{problema}[algoritmo digital]\label{prob:pn-digital}
Con el mismo PID: (a) escriba el algoritmo posicional y el incremental con $T_s=0.1$\,s; (b) deduzca el incremental del posicional; (c) \etiquetaM{} compruebe que dan la misma salida con $\Gp=1/(10s+1)$; (d) \etiquetaM{} repita con $T_s=0.01$, $0.5$ y $1$\,s y halle el mayor $T_s$ con el que $\OS$ y $\ts$ empeoran menos del $5\%$ frente al continuo; (e) \etiquetaEstrella{} cuantice el error a 12 bits en $\pm10$ y mida el ruido en $u_c$ para cada $T_s$; (f) agregue una saturación y explique por qué el incremental sufre menos \emph{wind-up}.
\end{problema}

\begin{problema}[\emph{kicks}]\label{prob:pn-kicks}
Con $\Gp=1/(s^2+10s+20)$ y el PID paralelo $350/300/50$, compare las estructuras estándar, PI-D e I-PD. (a) Escriba $U_c(s)$ en función de $R$ y $Y$ para cada una. (b) Con el teorema del valor inicial, halle $u_c(0^+)$ ante un escalón de referencia; muestre que en la estándar no está acotado. (c) Repita con la derivada filtrada de constante $\tau_f$ y evalúe para $\tau_f=0.01$, $0.05$ y $0.1$\,s. (d) \etiquetaM{} Simule y tabule $|u_c|_{\max}$, $\OS$ y $\ts$. (e) \etiquetaM{} Agregue una saturación en $[-100,\,100]$ y diga cuáles saturan y cuánto tiempo. (f) Muestre que PI-D e I-PD son casos de la ponderación con $\beta$ y $\gamma$.
\end{problema}

\begin{problema}[\emph{wind-up}]\label{prob:pn-windup}
Con el ejemplo~\ref{ej:pn-windup}: (a) \etiquetaM{} simule sin saturación y reporte $\OS$, $\ts$ y el pico de $u_c$; (b) agregue la saturación y reporte cuánto tiempo satura y cuánto empeoran $\OS$ y $\ts$; (c) grafique el término integral y marque el instante en que el error cambia de signo y el instante en que el actuador sale de saturación; explique la diferencia; (d) implemente el recálculo con $T_t=\sqrt{\Ti\Td}$ y $T_t=\Ti$; (e) recorra $T_t$ entre $0.5$ y $20$\,s y grafique $\OS$, $\ts$ e IAE; (f) implemente la integración condicional y compárela con el mejor recálculo.
\end{problema}

\begin{problema}[dos grados de libertad]\label{prob:pn-beta}
Con el ejemplo~\ref{ej:pn-beta}: (a) verifique por Routh que el PI agresivo es estable; (b) \etiquetaM{} compare el rechazo de un escalón de perturbación a la entrada con el PI conservador ($\Kc=2$, $\Ti=1$); (c) simule la respuesta a la referencia y constate el conflicto; (d) obtenga $G_f(s)$ para un PI con ponderación $\beta$; (e) muestre que $G_f$ no aparece en la función de transferencia de la perturbación; (f) tabule $\OS$ y $\ts$ para $\beta=0$, $0.3$, $0.5$ y $1$; ¿qué valores de $\beta$ dan $\OS\le30\%$? Explique por qué ningún $\beta$ baja de unos $27\%$.
\end{problema}

%% ---------------------------------------------------------------------
\section{Soluciones seleccionadas}
%% ---------------------------------------------------------------------

\paragraph{Problema~\ref{prob:pn-formas}}
(b) $\Kc^I=2.5$, $\Ti^I=10$\,s, $\Td^I=1.6$\,s. (e) $\Ti=4<8$; ceros $-0.25\pm j0.25$. (f) $\Kc=350$, $\Ti=1.167$\,s, $\Td=0.143$\,s; como $\Ti\ge4\Td$, sí: $\Kc^S=300$, $\Ti^S=1$\,s, $\Td^S=0.167$\,s (ceros en $-1$ y $-6$).

\paragraph{Problema~\ref{prob:pn-mra}}
(a) Valor final $0.05$, $\zeta=1.12$, $\wn=4.47$\,rad/s, $\tr=0.88$\,s, $\ts=1.60$\,s. (b) $\ess=\frac{20}{20+\Kc}$, $\wn=\sqrt{20+\Kc}$, $\zeta=\frac{5}{\sqrt{20+\Kc}}$; para $\Kc=10$, $100$, $300$: $\ess=0.667$, $0.167$, $0.0625$; $\OS=0.1\%$, $20\%$, $40\%$. (c) Cero en $-30$ y $-15$; $\OS=15.3\%$ y $5.8\%$ frente a $12.0\%$ y $0.8\%$ sin el cero. (d) $0<K_i<3200$; $\OS=33\%$ ($\ts=4.7$\,s) y $36\%$ ($\ts=1.2$\,s). (e) Aquí, subir $\Kc$ de $10$ a $100$ \emph{baja} $\ts$ (de $0.88$ a $0.76$\,s) en lugar de subirlo; de $100$ a $300$ casi no lo cambia. Las tendencias de la tabla no son leyes.

\paragraph{Problema~\ref{prob:pn-banda}}
(a) $\Kc=4$ (adimensional) y $\Kc=2$\,\%/$\grados$C. (b) $e=-0.83\,\grados$C (el \emph{bias} de $50\%$ llevaría la PV a $125\,\grados$C). (c) $u_0=48\%$. (d) $+2.67$ y $-2.18\,\grados$C. (f) $\Kc\ge4.6$\,\%/$\grados$C, es decir, $PB\le10.9\%$.

\paragraph{Problema~\ref{prob:pn-ruido}}
(a) $0.006$, $0.6$ y $6$. (b) $\Kc N$ en la derivada, $\Kc(1+N)$ en el controlador. (c) A partir de $N/\Td$: $2.5$, $5$ y $10$\,rad/s; amplitudes $0.15$, $0.30$ y $0.60$.

\paragraph{Problema~\ref{prob:pn-digital}}
(a) $u_c[k]=3\bigl(e[k]+0.0125\sum e[i]+20(e[k]-e[k-1])\bigr)$ y $u_c[k]=u_c[k-1]+63.04\,e[k]-123\,e[k-1]+60\,e[k-2]$.

\paragraph{Problema~\ref{prob:pn-kicks}}
(b) $u_c(0^+)=\Kc+K_d\,\delta$: un impulso. (c) $u_c(0^+)=\Kc+K_d/\tau_f=5350$, $1350$ y $850$. (d) $|u_c|_{\max}\approx5300$, $350$ y $20$; $\OS=0$, $5.3$ y $0\%$; $\ts=0.84$, $2.06$ y $4.29$\,s. (f) PI-D: $\beta=1$, $\gamma=0$; I-PD: $\beta=\gamma=0$.

\paragraph{Problema~\ref{prob:pn-windup}}
(a) Sin saturación, $\OS=67\%$ y el pico de $u_c$ es de $32$ (el \emph{derivative kick} con $N=10$). (b) $25$\,s saturado, $\OS=38\%$, $\ts=45$\,s. (d) $T_t=4$: $\OS=3.3\%$, $\ts=30$\,s; $T_t=8$: $\OS=11\%$, $\ts=32$\,s. (f) $\OS=3.7\%$, $\ts=29$\,s, sin parámetro que sintonizar.

\paragraph{Problema~\ref{prob:pn-beta}}
(a) $s^3+6s^2+25s+66.7$: $6\cdot25=150>66.7$. (b) Pico $0.032$ frente a $0.114$. (f) $\OS=27\%$, $29\%$, $33\%$ y $51\%$; $\OS\le30\%$ con $\beta\lesssim0.35$.

\respuestaspausas

